\section{Conclusion}\label{sec:conclusion}
RAMSES addresses the practical difficulty of coordinating VRAM, DRAM, and NVMe in a single inference node. The revised model distinguishes resident demand from transferred bytes and represents asynchronous overlap explicitly; the controller turns measured pressure into documented admission, prefetch, and eviction actions with hysteresis. On the reported laboratory testbed, RAMSES reduces inference latency by 34--41\% across five models and cuts whole-node energy per token by 47\% (energy--delay product by 70\%) while raising token throughput by 61\% relative to PyTorch; model-load time and peak VRAM fall by 60\% and 15\%. On a named industrial defect-inspection task RAMSES preserves baseline accuracy (output equivalence above 0.999). A 72-hour parameterized synthetic trace further tests sustained-load behavior.

These findings are empirical and configuration-dependent. They do not establish a worst-case latency, TSN/PLC compatibility, closed-loop stability, certified absolute (rack-level) energy accounting, or IEC~61508 compliance. The main result is therefore a reproducible single-node systems result: explicit cross-tier coordination can help in measured coordination-dominated operating points, whereas capacity and transfer pressure bound its benefit. More broadly, RAMSES contributes a falsifiable method---measure pressure, select a regime-specific action, and test the predicted boundary against a controller-disabled counterfactual---rather than presenting offload as a universal optimization. Field traces, certified rack-level power instrumentation, multi-node failure analysis, and hardware- or software-in-the-loop validation are required before stronger industrial deployment claims can be made.
